\subsection{仿紧空间}

我们回顾（第 I 章，§ 9，第 10 号）：拓扑空间 X 称为仿紧的，如果它是豪斯多夫的且 X 的每个开覆盖都有局部有限的开加细。

命题 4. 每个仿紧空间都是正规的。

这是下述引理的推论：

引理 2. 设 A、B 是仿紧空间 X 的两个不相交闭子集。若对每个 \(x \in \mathbf{A}\) 存在开邻域 \(\mathbf{V}_x\) （\(x\) 的）与 B 的开邻域 \(\mathbf{W}_x\) ，二者不相交，则存在 A 的开邻域 T 与 B 的开邻域 U，二者不相交。

暂且假定这个引理成立，我们可把它应用于 B 由单点组成的情形，因为 X 是豪斯多夫的，由此得出 X 是正则的。然后又可把引理 2 应用于 X 的任意两个不相交闭子集，这就表明公理 ( \(O_{V}^{\prime}\) ) 满足。

为证引理，考虑由 \(\mathbb{C}\mathbf{A}\) 与诸 \(\mathbf{V}_x\) （其中 \(x \in \mathbf{A}\) ）组成的 X 的开覆盖；设 \((\mathbf{T}_t)_{t \in \mathbf{I}}\) 是这个覆盖的局部有限开加细。于是按定义，若 \(\mathbf{A} \cap \mathbf{T}_t \neq \emptyset\) ，则存在 \(x_t \in \mathbf{A}\) 使 \(\mathbf{T}_t \subset \mathbf{V}_{x_t}\) 。设 T 是诸 \(\mathbf{T}_t\) （与 A 相交者）之并的开集，我们来证明存在 B 的开邻域 U 不与 T 相交。对每个 \(y \in \mathbf{B}\) ，有开邻域 \(\mathbf{S}_y\) （\(y\) 的）只与有限个集 \(\mathbf{T}_t\) 相交；设 J 是 I 中由使 \(\mathbf{T}_t\) 与 \(\mathbf{S}_y\) 及 A 都相交的那些指标 t 组成的有限子集；若令 \(\mathbf{U}_y = \mathbf{S}_y \cap \bigcap_{t \in \mathbf{J}} \mathbf{W}_{x_t}\) ，则 \(\mathbf{U}_y\) 是 \(y\) 的开邻域，不与任何和 A 相交的 \(\mathbf{T}_t\) 相交，因而 \(\mathbf{U}_y \cap \mathbf{T} = \emptyset\) 。令 \(\mathbf{U} = \bigcup_{y \in \mathbf{B}} \mathbf{U}_y\) ；则 \(\mathbf{U}'\) 是 B 的不与 T 相交的开邻域，引理得证。

存在不是仿紧的正规空间（习题 19）。

推论 1. 给定任一开覆盖 \((\mathbf{A}_t)_{t \in \mathbf{I}}\) （仿紧空间 \(\mathbf{X}\) 的），存在连续单位分解 \((f_t)_{t \in \mathbf{I}}\) （\(x\) 上的），从属于覆盖 \((\mathbf{A}_t)\) 。

设 \((\mathbf{U}_{\lambda})_{\lambda \in \mathbf{L}}\) 是 \(\mathbf{X}\) 的局部有限开覆盖，它加细覆盖 \((\mathbf{A}_{\iota})_{\iota \in \mathbf{I}}\) ；于是有映射 \(\varphi : \mathbf{L} \to \mathbf{I}\) 使 \(\mathbf{U}_{\lambda} \subset \mathbf{A}_{\varphi(\lambda)}\) 对每个 \(\lambda \in \mathbf{L}\) 成立。由命题 3 与 4，存在连续单位分解 \((g_{\lambda})_{\lambda \in \mathbf{L}}\) 从属于 \((\mathbf{U}_{\lambda})\) 。对每个 \(\iota \in \mathbf{I}\) ，令

\[
f _ {t} = \sum_ {\varphi (\lambda) = t} g _ {\lambda};
\]

这个和有定义且连续，因为诸 \(g_{\lambda}\) 的支集构成局部有限覆盖；此外，支集之并 \(\mathbf{B}_{t}\) ——由诸 \(g_{\lambda}\) （满足 \(\varphi(\lambda)=\iota\) 者）的支集组成——是闭集（第 I 章，§ 1，第 6 号，命题 4），且含于 \(\mathbf{A}_{t}\) 。由于 \(f_{t}(x)=0\) （每当 \(x\in\complement\mathbf{B}_{t}\) ），故 \(f_{t}\) 的支集含于 \(\mathbf{B}_{t}\) ，从而含于 \(\mathbf{A}_{t}\) 。另一方面，族 \(\mathbf{B}_{t}\) 是局部有限的，因为对每个 \(x\in\mathbf{X}\) 有 \(x\) 的邻域 V 与 L 的有限子集 H，使 V∩U \(_{\lambda}=\emptyset\) 对一切 \(\lambda\notin\mathbf{H}\) 成立，由此推出 V∩B \(_{t}=\emptyset\) 对一切 \(\iota\notin\varphi(\mathbf{H})\) 成立。最后，对每个 \(x\in\mathbf{X}\) 我们有

\[
\mathrm{I} = \sum_ {\lambda \in \mathbf {L}} g _ {\lambda} (x) = \sum_ {t \in \mathbf {I}} \left(\sum_ {\varphi (\lambda) = t} g _ {\lambda} (x)\right) = \sum_ {t \in \mathbf {I}} f _ {t} (x),
\]

于是证明完成。

推论 2. 若 F 是仿紧空间 X 的闭子集，则 F 在 X 中的每个邻域都包含 F 的一个闭（从而仿紧的）邻域。

由第 I 章，§ 9，第 10 号的命题 16，X 的任何闭子空间都是仿紧的；因此这个推论由命题 4 与公理 (O \(_{V}\) '') 即得。
% semantic-fragments: legacy-input-seam
\relax{}
